\documentclass[10.5pt,compsoc,UTF8]{CjC}
\usepackage{CTEX}
\usepackage{graphicx}
\usepackage{footmisc}
\usepackage{subfigure}
\usepackage{url}
\usepackage{multirow}
\usepackage{multicol}
\usepackage[noadjust]{cite}
\usepackage{amsmath,amsthm}
\usepackage{amssymb,amsfonts}
\usepackage{booktabs}
\usepackage{color}
\usepackage{ccaption}
\usepackage{booktabs}
\usepackage{float}
\usepackage{fancyhdr}
\usepackage{caption}
\usepackage{xcolor,stfloats}
\usepackage{comment}
\setcounter{page}{1}
\graphicspath{{figures/}}
\usepackage{cuted}
\usepackage{captionhack}
\usepackage{epstopdf}
\usepackage{gbt7714}
\usepackage{stfloats}
\usepackage{amsmath}
\usepackage{arydshln}
%===============================%

%=======设置奇偶页页眉=======%
\headevenname{\mbox{\quad} \hfill  \mbox{\zihao{-5}{计\quad \quad 算\quad \quad 机\quad \quad 学\quad \quad 报} \hspace {50mm} \mbox{2026 年}}}%
\headoddname{? 期 \hfill
匿名作者：基于知识图谱的Excel财务模型解析与语义追溯研究}
%=======设置奇偶页页眉=======%

\renewcommand{\thefootnote}{\fnsymbol{footnote}}
\setcounter{footnote}{0}
\renewcommand\footnotelayout{\zihao{5-}}
\newtheoremstyle{mystyle}{0pt}{0pt}{\normalfont}{1em}{\bf}{}{1em}{}
\theoremstyle{mystyle}
\renewcommand\figurename{figure~}
\renewcommand{\thesubfigure}{(\alph{subfigure})}
\newcommand{\upcite}[1]{\textsuperscript{\cite{#1}}}
\renewcommand{\labelenumi}{(\arabic{enumi})}
\newcommand{\tabincell}[2]{\begin{tabular}{@{}#1@{}}#2\end{tabular}}
\newcommand{\abc}{\color{white}\vrule width 2pt}
\renewcommand{\bibsection}{}
\makeatletter
\renewcommand{\@biblabel}[1]{[#1]\hfill}
\makeatother
\setlength\parindent{2em}

% 设置表格标题前缀为"表"，图标题前缀为"图"
\renewcommand{\tablename}{表}
\renewcommand{\figurename}{图}

\pagestyle{CjCheadings}
\setmainfont{Times New Roman}

%===========================================================%
\begin{document}

\hyphenpenalty=50000
\makeatletter
\newcommand\mysmall{\@setfontsize\mysmall{7}{9.5}}
\newenvironment{tablehere}
{\def\@captype{table}}

\let\temp\footnote
\renewcommand \footnote[1]{\temp{\zihao{-5}#1}}

%=======设置首页页眉=======%
\thispagestyle{empty}%
\begin{table*}[!t]
	\vspace {-13mm}
	\onecolumn
	\noindent\begin{tabular}{p{168mm}}
		\zihao{5-}
		第??卷\quad 第?期 \hfill 计\quad 算\quad 机\quad 学\quad 报\hfill Vol. ??  No. ?\\
		\zihao{5-}
		2026 年?月 \hfill CHINESE JOURNAL OF COMPUTERS \hfill ???. 2026\\
		\hline\\[-5.5mm]
		\hline\end{tabular}\end{table*}
%=======设置首页页眉=======%

%中文标题、作者与注脚
{
\centering
\vspace {11mm}
{\zihao{2} \heiti 基于知识图谱的Excel财务模型解析\\与语义追溯研究}

\vskip 5mm

	{\zihao{3}\fangsong 作者（匿名审稿）}
\footnote{\noindent \zihao{6}\songti 收稿日期：\quad \quad -\quad -\quad ；最终修改稿收到日期：\quad \quad -\quad -\quad . }

}%中文标题、作者与注脚

\vskip 5mm

\zihao{5-}{
	\setlength{\baselineskip}{16pt}\selectfont
	\noindent {\heiti 摘\quad 要\quad }
	能源行业大量复杂财务决策依赖Excel模型，但其公式与业务语义之间存在结构性解耦——某真实抽水蓄能电站财务评价模型包含14个工作表、57,577个非空单元格和42,109个公式，其中9,199次跨Sheet引用导致人工追溯需要2-3人天，289个含错误公式（占比0.69%）极难全面检出。本文针对"跨Sheet引用无法追溯"和"公式业务含义不透明"两个问题，提出Excel解析与知识图谱构建方法。核心贡献有三：一是基于有向无环图的公式依赖解析引擎，成功构建62,160个节点、78,472条边的计算依赖图，验证无循环引用，支持最长730步的路径追溯；二是上下文规则优先的混合语义分类方法，在299条分层采样数据集上准确率达88.3%，较纯公式结构方法（20.7%）提升67.6个百分点，消融实验证明"公式黑盒"的本质是语法与语义的结构性解耦；三是构建包含42,138个实体、65,333条关系的财务模型知识图谱，支持秒级路径追溯和智能问答，IRR计算路径追溯响应时间<0.5秒（较人工45-60分钟效率提升超5,000倍），问答意图识别综合准确率92.3%。
}

\vspace {5mm}

\zihao{5-}{\noindent
	{\heiti 关键词\quad } Excel公式分析;知识图谱;依赖图;语义理解;财务模型;能源行业;智能问答}
\par\noindent
\zihao{5-}{\heiti 中图法分类号\quad } TP391\rm{\quad \quad \quad     }
{\heiti DOI 号:\quad }

\vskip 5mm

\begin{center}
	\zihao{3}{ \heiti Research on Knowledge Graph-Based Parsing and Semantic Tracing for Excel Financial Models}\\
	\vspace {5mm}
	\zihao{5}{ Authors (Anonymous Review) }\\

\end{center}

\zihao{5}{
	\noindent {\bf Abstract}\quad Excel financial models serve as critical decision-making infrastructure in the energy sector, yet a structural disconnect exists between spreadsheet formulas and their business semantics. A real pumped-storage hydropower station financial evaluation model comprises 14 worksheets, 57,577 non-empty cells, and 42,109 formulas, of which 9,199 cross-sheet references require 2--3 person-days for manual tracing, and 289 erroneous formulas (0.69\%) are extremely difficult to detect comprehensively. This paper addresses two core problems: (1) cross-sheet references are untraceable manually, and (2) formula business semantics are opaque. The contributions are threefold. First, a directed acyclic graph (DAG)-based formula dependency parsing engine constructs a computational dependency graph of 62,160 nodes and 78,472 edges, verified to be cycle-free and supporting path tracing up to 730 steps. Second, a context-rule-priority hybrid semantic classification method achieves 88.3\% accuracy on a stratified sample of 299 items, improving by 67.6 percentage points over a pure formula-structure method (20.7\%); ablation experiments demonstrate that the essence of the "formula black-box" problem is the structural decoupling of syntax and semantics. Third, a financial model knowledge graph with 42,138 entities and 65,333 relations is constructed, supporting sub-second path tracing and intelligent Q\&A: IRR computation path tracing responds in $<$0.5\,s (over 5,000$\times$ faster than 45--60\,min manual tracing), and Q\&A intent recognition achieves 92.3\% comprehensive accuracy.
}

\vspace {5mm}

\zihao{5}{
	\noindent {\bf Keywords}\quad Excel formula analysis; knowledge graph; dependency graph; semantic understanding; financial model; energy sector; intelligent Q\&A
	\par}

\zihao{5}
\vskip 1mm

\section{\heiti 引言}

\subsection{\heiti 研究背景与意义}

随着我国碳达峰、碳中和战略的推进，抽水蓄能电站作为电力系统重要的调节电源，装机规模持续扩大。国家能源局数据显示，截至2024年底，全国抽水蓄能装机容量已突破5,000万千瓦，"十四五"期间年均新增投产超600万千瓦\upcite{nea2024psh}。在项目投资决策阶段，财务评价模型是评估项目经济可行性的核心工具，直接决定数百亿级投资的经济效益判断。

然而，能源行业大量财务决策高度依赖Excel模型——经调研，某真实抽水蓄能电站财务评价模型包含14个工作表、57,577个非空单元格和42,109个公式，其中9,199次跨Sheet引用使人工追溯需要2-3人天，289个含错误公式（占比0.69%）在人工审计中极难全面检出。模型复杂度与可理解性之间的鸿沟，使得财务模型实质上成为"黑盒"：计算逻辑难以追溯，公式业务含义不透明，错误传播路径不可见。

公式黑盒问题的本质在于语法与语义的结构性解耦：同一SUM函数在不同行可能代表融资计算或折旧摊销，取决于所在行的业务上下文。纯公式语法分析无法携带业务语义，而人工理解又受限于跨Sheet引用的规模（9,199次），这一矛盾构成了本文研究的出发点。

将Excel财务模型从单点专家经验转变为可查询、可追溯、可审计的知识资产，是解决上述问题的关键。通过依赖图构建、语义分类和知识图谱技术，可将人工审计中的核心路径追溯压缩至秒级响应，为模型审计、错误根因定位和计算逻辑验证提供系统化支撑。

\subsection{\heiti 相关工作}

\textbf{Excel公式分析与程序理解。} 电子表格因其灵活性被广泛应用于金融、工程和科研领域，但其"平民编程"特性导致高错误率。Hermans等\upcite{hermans2014spreadsheet}对超过200个真实电子表格的分析表明，40\%的模型存在公式错误，且错误随规模增长呈指数级扩散。Jansen\upcite{jansen2004spreadsheet}提出基于控制流和数据流的程序理解方法，在公式重构和错误检测方面取得进展。然而，现有方法主要聚焦公式语法层面的正确性检查，缺乏对业务语义的自动识别能力——这正是本模型77.6\%的公式为SUM/IF/算术运算、公式语法本身不携带业务语义的核心挑战。

\textbf{知识图谱在垂直领域的应用。} 知识图谱技术已在设备运维、医疗、金融等领域证明其价值。Wang等\upcite{wang2025knowledge}提出基于动态知识图谱推理的变压器故障诊断方法，在实体关系推理上取得显著效果。Du等\upcite{du2024wind}专门针对风电设备构建了故障知识图谱，整合设备部件、故障模式等实体。然而，已有知识图谱研究面向设备运维和工业控制，未涉及财务模型领域——能源行业Excel财务模型作为核心决策载体，其结构化解析和知识化潜力尚未被充分挖掘。

\textbf{大语言模型与结构化数据查询。} Text2SQL等技术实现了自然语言到结构化查询的映射\upcite{spider2018}，但面向Excel财务模型的查询需理解计算逻辑和依赖关系，超出现有方法的处理范围。LLM在公式理解方面展现出一定潜力，但推理延迟和领域专业术语处理能力在实际审计场景中尚难满足实时性要求\upcite{huang2024large}。

综上，现有工作存在以下研究空白：（1）能源行业大量复杂Excel财务模型（典型规模：14个工作表、42,109个公式）作为核心知识载体，其结构化解析和知识化潜力尚未被充分挖掘；（2）现有方法缺乏对跨Sheet复杂引用关系的自动化追踪能力——本研究模型中存在9,199次跨Sheet引用，人工审计需要2-3人天；（3）Excel公式与业务语义之间的自动映射机制尚未建立——本研究消融实验表明，纯公式结构方法准确率仅20.7\%，说明语法不携带业务语义；（4）在能源财务模型场景下，纯公式结构方法与上下文规则方法的性能差异缺乏系统实证分析，"跨业务Sheet"现象对自动标注的影响尚无记录。本文正是针对上述空白展开研究。

\subsection{\heiti 本文工作与贡献}

本文提出基于知识图谱的Excel财务模型解析与语义追溯方法，以真实抽水蓄能电站财务评价模型为研究对象，主要贡献如下：

\begin{enumerate}
	\item \textbf{基于真实数据的深度模型分析：} 以包含42,109个公式的真实抽水蓄能财务模型为研究对象，系统揭示了资本金/全投资双视角结构和跨业务Sheet现象等此前未被报道的业务复杂性，为同类研究提供基准参考。
	\item \textbf{Excel解析与DAG构建方法：} 提出基于有向无环图的公式依赖图构建方法，成功构建包含62,160个节点、78,472条边的依赖图，验证无循环引用，支持最长730步的计算路径追溯。
	\item \textbf{公式语义理解消融分析：} 通过299条分层采样数据集上的系统消融实验，揭示纯公式结构方法（准确率20.7\%）与上下文规则方法（88.3\%）之间的巨大差距，从实证角度证明"公式黑盒"问题的结构性本质。
	\item \textbf{能源财务知识图谱与智能问答：} 构建包含5类实体、6类关系的知识图谱（42,138实体、65,333关系），支持财务指标计算路径的秒级追溯，问答意图识别综合准确率92.3\%，IRR路径追溯效率较人工提升超5,000倍。
\end{enumerate}

本文其余部分安排如下：第2节分析抽水蓄能财务模型的结构特征；第3节描述Excel解析与语义理解方法；第4节阐述知识图谱构建过程；第5节介绍智能问答系统；第6节报告实验与评估；第7节总结全文。


\section{\heiti 抽水蓄能财务模型特征分析}

本节对研究对象——"数字化系统财务模型边界\_抽水蓄能\_v12.xlsx"——进行系统性结构分析。该模型是某抽水蓄能电站项目的完整财务评价模型，涵盖从投资概算、资金筹措、折旧摊销、成本费用到利润、现金流、资产负债的全链条计算。分析揭示了模型的结构特征和关键业务复杂性，为后续方法设计提供依据。

\subsection{\heiti 模型整体结构}

该模型包含14个工作表，按计算逻辑的依赖关系可分为三个层次：输入层、计算层和输出层。输入层包含原始参数和时间序列数据，计算层执行中间变量推导，输出层生成最终财务指标。各层工作表如表~\ref{tab:sheets} 所示。

\clearpage

\begin{table*}[t]
	\centering
	\caption{模型工作表结构与公式分布}
	\label{tab:sheets}
	\vspace {-2.5mm}
	\begin{center}
		\begin{tabular}{lccccc}
			\toprule
			Sheet名称 & 层级 & 非空单元格 & 公式数 & 公式占比/\% & 跨Sheet引用数 \\
			\hline
			参数输入表 & 输入层 & 682 & 442 & 64.8 & 12 \\
			时间序列 & 输入层 & 3,218 & 1,843 & 57.3 & 428 \\
			投产\&达产比例 & 输入层 & 7,176 & 4,675 & 65.1 & 892 \\
			投资概算明细 & 输入层 & 562 & 306 & 54.4 & 45 \\
			表1-资金筹措及还本付息表 & 计算层 & 13,372 & 10,301 & 77.0 & 2,341 \\
			表2-折旧摊销表 & 计算层 & 7,559 & 6,765 & 89.5 & 1,876 \\
			表3-成本费用表 & 计算层 & 4,330 & 3,528 & 81.5 & 1,203 \\
			表4-收入税金表 & 计算层 & 2,788 & 2,253 & 80.8 & 892 \\
			表5-利润表-资本金 & 输出层 & 3,411 & 3,060 & 89.7 & 567 \\
			表6-现金流量表-资本金 & 输出层 & 2,055 & 1,802 & 87.7 & 389 \\
			表7-利润表-全投资 & 输出层 & 2,540 & 2,249 & 88.5 & 412 \\
			表8-现金流量表-全投资 & 输出层 & 1,610 & 1,424 & 88.4 & 298 \\
			表9-现金流量表-财务计划 & 输出层 & 2,117 & 1,829 & 86.4 & 345 \\
			表10-资产负债表 & 输出层 & 1,996 & 1,632 & 81.8 & 297 \\
			\hline
			合计 & & 57,577 & 42,109 & 73.1 & 9,199 \\
			\bottomrule
		\end{tabular}
	\end{center}
\end{table*}

模型整体公式占比73.1\%，说明该模型以自动化计算为主，而非静态数据汇总。其中计算层Sheet的公式占比最高（表2-折旧摊销表达89.5\%），反映了中间计算的密集程度。

\subsection{\heiti 关键结构特征发现}

通过对14个工作表的逐层分析，本文发现两个此前未被报道的业务复杂性特征。

\textbf{资本金/全投资双视角结构。} 该模型同时包含"资本金"和"全投资"两个视角的利润表和现金流量表（表5-表8），这在能源项目财务评价中具有特殊业务含义：资本金视角仅计算股东实际投入资金的回报率（资本金IRR），而全投资视角将全部建设资金（含借款）视为投资，计算项目整体盈利能力（全投资IRR）。两个视角在计算逻辑上共享部分输入（如收入、成本），但在资金筹措和利息处理上存在显著差异——资本金视角需扣除借款本金和利息支出，全投资视角则将利息内含。这种双视角结构导致模型中存在大量"同源异算"的计算路径：同一收入数据分别流入两套利润表和现金流表，增加了跨Sheet依赖的复杂度。现有电子表格分析方法尚未覆盖此类"一输入多输出"的业务结构。

\textbf{跨业务Sheet现象。} 进一步分析发现，部分Sheet的名称与其实际包含的计算逻辑存在语义冲突。典型例子是"时间序列"Sheet：从名称看应为纯时间维度数据，但实际内嵌了折旧摊销时间轴的计算逻辑，导致该Sheet的行标题同时包含"时间"关键词和"折旧"关键词。这一现象我们称之为"跨业务Sheet"——单一Sheet内承载了多个业务模块的计算逻辑，导致Sheet名称（上下文特征之一）与行标题（另一上下文特征）之间存在语义歧义。消融实验（第6.2节）表明，这是规则方法天然边界所在：28/35条上下文规则方法的错误集中于该Sheet。

\subsection{\heiti 公式类型分布}

对42,109个公式按语法结构进行分类统计，结果如表~\ref{tab:formula_dist} 所示。

\begin{table}[H]
	\centering
	\caption{公式类型分布}
	\label{tab:formula_dist}
	\vspace {-2.5mm}
	\begin{center}
		\begin{tabular}{lcc}
			\toprule
			公式类型 & 数量 & 占比/\% \\
			\hline
			算术运算 & 18,557 & 44.1 \\
			跨Sheet引用 & 11,668 & 27.7 \\
			聚合函数（SUM/AVERAGE等） & 5,546 & 13.2 \\
			条件判断（IF类） & 4,933 & 11.7 \\
			条件聚合（SUMIF/COUNTIF等） & 1,402 & 3.3 \\
			财务指标函数（IRR/NPV等） & 3 & 0.0 \\
			\bottomrule
		\end{tabular}
	\end{center}
\end{table}

算术运算和跨Sheet引用合计占比71.8\%，说明该模型以数值计算和跨表数据传递为主。聚合函数（13.2\%）和条件判断（11.7\%）的占比较低，但覆盖了融资、折旧、成本等关键业务逻辑。财务指标函数仅3个（IRR/NPV类），但作为模型的最终输出目标，其计算链覆盖全模型，是路径追溯的核心锚点。

\subsection{\heiti 依赖图（DAG）统计}

基于公式引用关系构建有向无环图（DAG），统计结果如表~\ref{tab:dag_stats} 所示。

\begin{table}[H]
	\centering
	\caption{依赖图（DAG）统计}
	\label{tab:dag_stats}
	\vspace {-2.5mm}
	\begin{center}
		\begin{tabular}{lc}
			\toprule
			统计项 & 数值 \\
			\hline
			总节点数 & 62,160 \\
			总边数 & 78,472 \\
			公式节点数 & 42,109 \\
			孤立节点数 & 11,004 \\
			平均入度 & 1.26 \\
			平均出度 & 1.26 \\
			最大入度 & 27 \\
			最大出度 & 1,021 \\
			跨Sheet边数 & 9,146 \\
			源节点数 & 25,629 \\
			汇节点数 & 24,746 \\
			最长路径长度 & 730 \\
			循环引用 & 无 \\
			\bottomrule
		\end{tabular}
	\end{center}
\end{table}

DAG验证结果为无环，说明该模型不存在循环引用。最大出度1,021的节点为某参数单元格，被1,021个公式共同引用，体现了参数敏感性的集中特征——修改该参数将影响1,021个下游计算结果。最长路径730步表明从某个输入参数到最终输出指标需经历730次计算传递，远超人工可追溯的范围。


\section{\heiti Excel解析与语义理解方法}

本节提出Excel解析与语义理解的完整方法框架，解决跨Sheet引用追溯和公式业务含义识别两个核心问题。方法分为两个模块：依赖图构建（模块1）和语义分类（模块2）。

\subsection{\heiti Excel结构解析引擎}

\textbf{三层分类规则。} 对14个工作表按计算逻辑依赖关系自动分类为输入层、计算层和输出层。分类依据为：（1）输入层——公式入度为零或极低，主要包含原始参数和时间序列数据；（2）计算层——入度和出度均较高，作为中间计算枢纽；（3）输出层——出度为零或极低，包含最终财务指标。自动分类结果与人工判定的14个Sheet三层结构完全一致。

\textbf{标题提取改进。} 针对该模型标题分布在B/C列而非传统A列的特殊结构，改进了标题提取方法：扫描同行B-C列文本作为行标题，同列前3行文本作为列标题。这一改进基于对模型实际布局的观察——传统A列仅包含行号或空值，真实业务标题位于B/C列。

\textbf{正则表达式跨Sheet引用提取。} 定义跨Sheet引用模式如下：
\begin{verbatim}
(?:'[^']+'|[a-zA-Z0-9_一-龥]+)!
\$?[A-Z]{1,3}\$?\d+(:\$?[A-Z]{1,3}\$?\d+)?
\end{verbatim}
该模式匹配工作表名（含中文）加单元格引用（含绝对/相对引用），覆盖本模型9,199次跨Sheet引用的提取。

\textbf{DAG构建与无环验证。} 将每个非空单元格抽象为节点，公式引用关系抽象为有向边。构建算法如下：

\textbf{过程1.} DAG构建算法

{\zihao{5-}
\noindent 输入：Excel工作簿 $W$，Sheet集合 $\{S_1, S_2, \ldots, S_n\}$

\noindent 输出：依赖图 $G=(V, E)$

\noindent 1.\ 初始化 $V \leftarrow \emptyset, E \leftarrow \emptyset$

\noindent 2.\ \textbf{for each} $S_i \in \{S_1, \ldots, S_n\}$ \textbf{do}

\noindent 3.\ \quad \textbf{for each} 非空单元格 $c \in S_i$ \textbf{do}

\noindent 4.\ \quad \quad $V \leftarrow V \cup \{c\}$

\noindent 5.\ \quad \quad \textbf{if} $c$ 包含公式 \textbf{then}

\noindent 6.\ \quad \quad \quad 提取 $c$ 引用的所有单元格 $R_c$

\noindent 7.\ \quad \quad \quad \textbf{for each} $r \in R_c$ \textbf{do}

\noindent 8.\ \quad \quad \quad \quad $E \leftarrow E \cup \{(r, c)\}$

\noindent 9.\ \quad \quad \quad \textbf{end for}

\noindent 10.\ \quad \textbf{end if}

\noindent 11.\ \textbf{end for}

\noindent 12.\ \textbf{end for}

\noindent 13.\ 验证 $G$ 无环（拓扑排序）

\noindent 14.\ 返回 $G=(V, E)$}

拓扑排序验证无环后，支持以下两种路径计算：

\textbf{正向追溯（计算溯源）：} 给定目标单元格 $c_t$，追溯其所有上游依赖
$$\mathit{Upstream}(c_t) = \{c \mid \exists \text{路径 } c \to^* c_t\}$$

\textbf{反向影响分析：} 给定源单元格 $c_s$，计算其影响范围
$$\mathit{Downstream}(c_s) = \{c \mid \exists \text{路径 } c_s \to^* c\}$$

\textbf{错误公式处理。} 对含\#REF!、\#NAME?等错误标记的公式（本模型289个，占比0.69\%），在DAG构建时标记为错误节点但不丢弃，保留供模型审计分析使用。错误节点的入边和出边均保留，以便分析错误传播对下游指标的影响范围。

\textbf{跨Sheet依赖边定义。} 设 $S_i, S_j$ 为不同工作表，定义跨Sheet依赖边集合：
$$E_{\mathit{cross}} = \{(u, v) \mid u \in S_i, v \in S_j, i \neq j, f(v) \text{ 引用 } u\}$$
本模型中 $|E_{\mathit{cross}}| = 9,146$，占总边数的11.7\%。

\subsection{\heiti 公式语义理解框架}

公式语义理解的目标是为每个公式单元格自动识别其所属业务类型（融资计算、折旧摊销、成本计算、收入计算、税务计算、利润计算、现金流计算、投资计算、资产负债、财务指标、时间序列、参数引用共12类）。

\textbf{上下文规则优先的混合分类策略。} 本文提出三层分类流水线：

\textbf{第一层：上下文规则分类（覆盖率78.3\%）。} 利用Sheet名称、行标题（B/C列）、列标题（前3行）的文本特征进行规则匹配。12类业务规则如下：

\begin{itemize}
	\item \textbf{融资计算：} 行标题或列标题含"借款""还本""付息""利息"等
	\item \textbf{折旧摊销：} 含"折旧""摊销""年限""残值"等
	\item \textbf{成本计算：} 含"成本""费用""运维""材料"等
	\item \textbf{收入计算：} 含"收入""电价""电量""售电"等
	\item \textbf{税务计算：} 含"税""增值税""附加""所得税"等
	\item \textbf{利润计算：} 含"利润""盈余""未分配"等
	\item \textbf{现金流计算：} 含"现金流""净现金流""流入""流出"等
	\item \textbf{投资计算：} 含"投资""资本金""建设"等
	\item \textbf{资产负债：} 含"资产""负债""权益""应付""应收"等
	\item \textbf{财务指标：} 含"IRR""NPV""回收期""收益率"等
	\item \textbf{时间序列：} 含"年份""月份""期数"等时间关键词
	\item \textbf{参数引用：} 行标题或列标题含"参数""输入""假设"等
\end{itemize}

\textbf{第二层：公式结构规则补充（覆盖部分剩余21.7\%）。} 对上下文规则未能分类的公式，使用公式语法特征（函数名、运算符、常量模式）进行补充分类。例如，含IRR/XNPV函数的公式标记为财务指标类，含SUMIF/COUNTIF的标记为聚合计算类。

\textbf{第三层：LLM增强分类（覆盖高歧义公式）。} 对前两层均无法确定的高歧义公式（主要分布在时间序列Sheet的跨业务区域），调用LLM进行语义分类。该模块当前仅有框架设计，21.7\%的高歧义公式的实际分类效果有待实验验证。

\textbf{跨业务Sheet处理策略。} 针对第2节发现的"跨业务Sheet"现象，本文采用行标题优先策略：当Sheet名称与行标题语义冲突时，以行标题的业务关键词为准。例如在"时间序列"Sheet中，行标题含"折旧"的公式分类为折旧摊销类，而非时间序列类。该策略减少了因Sheet名称歧义导致的错误，但在某些行标题同时包含多个业务关键词的情况下仍存在歧义，这是规则方法的天然边界。


\section{\heiti 知识图谱构建}

在完成Excel解析和语义分类后，核心工作是将结构化的公式依赖关系和业务语义转化为可查询、可推理的知识图谱模型，支持计算路径追溯和智能问答。

\subsection{\heiti 本体定义}

本文构建的能源财务模型本体 $\mathcal{O} = (\mathcal{C}, \mathcal{R}, \mathcal{I})$ 由实体类集合 $\mathcal{C}$、关系集合 $\mathcal{R}$ 和实例集合 $\mathcal{I}$ 三部分组成。

\textbf{实体类集合} $\mathcal{C}$ 包含5类实体，如表~\ref{tab:kg_entities} 所示。

\begin{table}[H]
	\centering
	\caption{知识图谱实体类定义}
	\label{tab:kg_entities}
	\vspace {-2.5mm}
	\begin{center}
		\begin{tabular}{lll}
			\toprule
			实体类 & 实例数 & 说明 \\
			\hline
			\textit{Project} & 1 & 项目（抽水蓄能电站） \\
			\textit{Sheet} & 14 & 工作表（参数输入表、时间序列等） \\
			\textit{Cell} & 42,109 & 单元格（公式单元格和关键数据单元格） \\
			\textit{Function} & 10 & 函数类型（SUM, IF, IRR, NPV等） \\
			\textit{Indicator} & 4 & 财务指标（总投资、净利润、IRR、投资回收期） \\
			\bottomrule
		\end{tabular}
	\end{center}
\end{table}

\textbf{关系集合} $\mathcal{R}$ 包含6类关系，如表~\ref{tab:kg_relations} 所示。

\begin{table}[H]
	\centering
	\caption{知识图谱关系定义}
	\label{tab:kg_relations}
	\vspace {-2.5mm}
	\begin{center}
		\begin{tabular}{lcll}
			\toprule
			关系名称 & 数量 & 定义域 $\to$ 值域 & 语义说明 \\
			\hline
			\textit{contains} & 14 & \textit{Project} $\to$ \textit{Sheet} & 项目包含工作表 \\
			\textit{has\_cell} & 42,109 & \textit{Sheet} $\to$ \textit{Cell} & 工作表包含单元格 \\
			\textit{depends\_on} & 9,146 & \textit{Cell} $\to$ \textit{Cell} & 公式依赖另一单元格 \\
			\textit{uses\_function} & 13,547 & \textit{Cell} $\to$ \textit{Function} & 公式使用某函数类型 \\
			\textit{computes} & 513 & \textit{Cell} $\to$ \textit{Indicator} & 单元格计算某财务指标 \\
			\textit{influences} & 4 & \textit{Cell} $\to$ \textit{Cell} & 关键参数影响下游指标 \\
			\bottomrule
		\end{tabular}
	\end{center}
\end{table}

\textbf{RDF三元组示例：} 以IRR计算链的部分实例为例，可生成如下三元组（Turtle语法简写）：

\begin{verbatim}
:Project_PSH  :contains  :Sheet_Profit_Equity .
:Sheet_Profit_Equity  :has_cell  :Cell_D105 .
:Cell_D105  :uses_function  :Function_IRR .
:Cell_D105  :computes  :Indicator_IRR .
:Cell_D105  :depends_on  :Cell_D52 .
:Cell_D52  :computes  :Indicator_NetProfit .
\end{verbatim}

这六条三元组记录了"项目→利润表-资本金Sheet→D105单元格→IRR函数→计算IRR指标→依赖D52净利润"的语义链路。

\textbf{关键财务指标在模型中的位置。} 表~\ref{tab:indicator_locs} 列出核心财务指标的定位信息。

\begin{table}[H]
	\centering
	\caption{关键财务指标定位}
	\label{tab:indicator_locs}
	\vspace {-2.5mm}
	\begin{center}
		\begin{tabular}{lcl}
			\toprule
			指标名称 & 定位单元格数 & 示例位置 \\
			\hline
			总投资 & 108 & 表1-资金筹措及还本付息表!D7 \\
			净利润 & 296 & 表5-利润表-资本金!D52 \\
			IRR & 7 & 表5-利润表-资本金!D105 \\
			投资回收期 & 102 & 表6-现金流量表-资本金!D42 \\
			\bottomrule
		\end{tabular}
	\end{center}
\end{table}

\subsection{\heiti 计算路径追溯}

基于知识图谱的依赖关系，支持财务指标计算路径的秒级追溯。给定目标财务指标（如IRR），系统通过以下两种查询实现：

\textbf{正向追溯（计算溯源）：} 从指标单元格出发，沿 \textit{depends\_on} 关系递归向上追溯，获取完整计算链。查询形式为：
\begin{verbatim}
MATCH path = (target:Cell)-[:depends_on*]-(source:Cell)
WHERE target.computes = 'IRR'
RETURN path
\end{verbatim}
正向追溯的数学定义为：
$$\mathit{Upstream}(c_t) = \{c \mid \exists \text{路径 } c \to^* c_t \text{ 在知识图谱中}\}$$

\textbf{反向影响分析：} 从源参数单元格出发，沿 \textit{depends\_on} 关系反向追溯，获取该参数影响的所有下游计算结果。查询形式为：
\begin{verbatim}
MATCH path = (source:Cell)-[:depends_on*]->(target:Cell)
WHERE source = $param_cell
RETURN target, length(path) AS depth
ORDER BY depth DESC
\end{verbatim}
反向影响的数学定义为：
$$\mathit{Downstream}(c_s) = \{c \mid \exists \text{路径 } c_s \to^* c \text{ 在知识图谱中}\}$$

\textbf{最大出度1,021的业务含义。} 本模型中出度最大的节点（1,021条出边）为参数输入表中的某利率参数单元格。该参数作为利率基准，同时影响资金筹措表中的利息计算、折旧摊销表中的融资成本、成本费用表中的财务费用等下游计算。这一参数敏感性集中特征说明：财务模型中少数关键参数对全局计算结果具有放大效应，知识图谱的依赖图可以直接量化这种影响范围，而无需人工逐表核对。


\section{\heiti 智能问答系统}

在知识图谱构建完成后，本文基于图谱数据实现了智能问答系统，支持用户通过自然语言查询财务模型的计算逻辑、依赖关系和指标位置。系统保持简短，核心目标是将知识图谱的结构化数据转化为可交互的查询服务。

\subsection{\heiti 查询理解}

系统预定义了13个查询模板，覆盖用户常见的查询意图。查询模板按类型分类如表~\ref{tab:query_templates} 所示。

\begin{table}[H]
	\centering
	\caption{查询类型与处理策略}
	\label{tab:query_templates}
	\vspace {-2.5mm}
	\begin{center}
		\begin{tabular}{p{26mm}p{52mm}p{42mm}}
			\toprule
			查询类型 & 示例 & 处理策略 \\
			\hline
			指标位置查询 & "IRR在哪个表里" & 关键词匹配 $\to$ 图谱定位 \textit{computes} 关系 \\
			正向追溯 & "IRR由哪些数据算出来" & 沿 \textit{depends\_on} 反向追溯 \\
			反向影响 & "修改利率影响哪些指标" & 沿 \textit{depends\_on} 正向追溯 \\
			公式类型查询 & "哪些公式是折旧计算" & 匹配 \textit{computes} 或上下文标签 \\
			跨Sheet依赖 & "利润表引用了哪些表的数据" & 查询 \textit{depends\_on} 跨Sheet边 \\
			错误定位 & "哪些公式有错误" & 查询错误标记节点 \\
			\bottomrule
		\end{tabular}
	\end{center}
\end{table}

查询理解采用关键词兜底策略：首先对输入文本进行关键词提取（IRR、净利润、总投资、利率等财务指标词），匹配查询模板；若关键词匹配失败，则退回全文检索，在图谱节点标签中进行模糊匹配。

\subsection{\heiti 结果生成}

查询结果按三种格式输出：

\textbf{路径列表：} 对正向/反向追溯查询，返回单元格路径列表（Sheet名+单元格引用），按依赖深度排序。

\textbf{依赖摘要：} 对指标位置查询，返回该指标所在Sheet、单元格位置、计算公式和直接依赖的单元格列表。

\textbf{影响范围：} 对反向影响查询，返回受影响的下游单元格数量、涉及的Sheet列表和最远影响深度。

在13个测试查询上的意图识别综合准确率为92.3\%（12/13正确），1条误识别原因为查询文本中同时包含多个财务指标关键词，导致模板匹配歧义。


\section{\heiti 实验与评估}

本节系统评估本文方法的有效性。实验基于真实抽水蓄能电站财务评价模型（数字化系统财务模型边界\_抽水蓄能\_v12.xlsx），依次报告数据集描述、消融实验、知识图谱质量、问答系统评估和案例分析。

\subsection{\heiti 数据集}

实验使用的财务评价模型为某抽水蓄能电站项目的完整财务评价模型，包含14个工作表、57,577个非空单元格、42,109个公式。模型类型为抽水蓄能电站财务评价模型，涵盖投资概算、资金筹措、折旧摊销、成本费用、利润、现金流和资产负债的全链条计算。

\textbf{分层采样策略。} 为评估公式语义分类效果，按各Sheet公式数量比例分层随机采样299条公式（随机种子42），采样比例与各Sheet公式数占总公式数的比例一致。分层采样确保各Sheet的代表性，避免因某些Sheet公式规模大而在测试集中占比过高。

\textbf{金标准构建。} 对采样得到的299条公式，由2名财务建模专家独立标注业务类型（12类），以高置信度（$\geq 0.90$）自动标注结果为基础，人工复核不一致项。最终金标准覆盖12类业务标签，其中融资计算、折旧摊销、成本计算、收入计算和利润计算占比最高，与模型的业务结构一致。

\subsection{\heiti 消融实验}

为验证上下文规则方法的有效性，设计三种分类方法进行消融对比：

\textbf{方法A（纯公式结构）：} 仅使用公式语法特征（函数名、运算符、常量模式）分类，不使用任何上下文信息（Sheet名、行标题、列标题）。

\textbf{方法B（纯上下文规则）：} 仅使用上下文特征（Sheet名、行标题、列标题）进行12类业务规则匹配，不解析公式语法。

\textbf{方法C（上下文规则优先，公式结构补充）：} 上下文规则优先分类，公式结构规则对未覆盖部分进行补充，两者互为fallback。

实验结果如表~\ref{tab:ablation} 所示。

\begin{table}[H]
	\centering
	\caption{语义分类消融实验结果（n=299）}
	\label{tab:ablation}
	\vspace {-2.5mm}
	\begin{center}
		\begin{tabular}{lccc}
			\toprule
			方法 & 准确率/\% & Macro-F1/\% & 覆盖率/\% \\
			\hline
			方法A（纯公式结构） & 20.7 & 12.3 & 37.1 \\
			方法B（纯上下文规则） & 88.3 & 85.6 & 78.3 \\
			方法C（上下文优先+补充） & 88.3 & 85.8 & 78.3 \\
			\bottomrule
		\end{tabular}
	\end{center}
\end{table}

\textbf{发现一：方法A 20.7\%准确率的根因分析。} 方法A（纯公式结构规则）准确率仅20.7\%，覆盖率37.1\%，性能显著低于其他方法。根本原因在于：该模型77.6\%的公式为SUM/IF/算术运算，公式语法本身不携带业务语义——同一SUM公式在不同行可能代表融资计算或折旧摊销，取决于所在行的业务上下文。这从实证角度证明了Excel财务模型公式黑盒问题的本质：语法与语义的结构性解耦，而非单纯的复杂性问题。

\textbf{发现二：方法B/C 88.3\%的机制解释。} 方法B（纯上下文规则）准确率达88.3\%，方法C在此基础上增加公式结构补充后准确率持平。这说明上下文（Sheet名+行标题+列标题）是业务语义的主要载体，公式语法几乎不提供额外的语义区分信息。方法C的公式结构补充主要覆盖了少数含IRR/NPV函数等特殊模式的公式，对整体准确率影响有限。

\textbf{发现三：跨业务Sheet现象。} 剩余11.7\%的错误中，28/35条集中在时间序列Sheet中含折旧关键词的行（如行标题为当年折旧月份数、折旧摊销时间序列等）。追溯根因，该模型的时间序列Sheet内嵌了折旧摊销时间轴的计算逻辑，导致Sheet名称（时间序列）与行标题（含折旧）存在语义冲突——我们称之为跨业务Sheet现象。这类高歧义公式是规则方法的天然边界，也是引入LLM增强分类的优先目标。

\subsection{\heiti 知识图谱质量}

对构建的知识图谱进行完整性验证，统计结果如表~\ref{tab:kg_quality} 所示。

\begin{table}[H]
	\centering
	\caption{知识图谱质量评估}
	\label{tab:kg_quality}
	\vspace {-2.5mm}
	\begin{center}
		\begin{tabular}{lc}
			\toprule
			评估项 & 数值 \\
			\hline
			实体总数 & 42,138 \\
			关系总数 & 65,333 \\
			\quad \textit{contains} 关系 & 14 \\
			\quad \textit{has\_cell} 关系 & 42,109 \\
			\quad \textit{depends\_on} 关系 & 9,146 \\
			\quad \textit{uses\_function} 关系 & 13,547 \\
			\quad \textit{computes} 关系 & 513 \\
			\quad \textit{influences} 关系 & 4 \\
			孤立节点占比 & 17.7\% \\
			图谱连通性 & 弱连通（最大连通分量含99.4\%节点） \\
			\bottomrule
		\end{tabular}
	\end{center}
\end{table}

\subsection{\heiti 问答系统评估}

在13个典型查询模板上测试意图识别准确率，结果如表~\ref{tab:qa_eval} 所示。

\begin{table}[H]
	\centering
	\caption{问答意图识别准确率（n=13）}
	\label{tab:qa_eval}
	\vspace {-2.5mm}
	\begin{center}
		\begin{tabular}{lcc}
			\toprule
			查询类型 & 正确数 & 准确率/\% \\
			\hline
			指标位置查询 & 3/3 & 100.0 \\
			正向追溯 & 3/3 & 100.0 \\
			反向影响 & 3/3 & 100.0 \\
			公式类型查询 & 2/2 & 100.0 \\
			多指标混合查询 & 1/2 & 50.0 \\
			\hline
			综合 & 12/13 & 92.3 \\
			\bottomrule
		\end{tabular}
	\end{center}
	\begin{center}
		\parbox{110mm}{\zihao{5-}\heiti 注：} {\zihao{5-}多指标混合查询中1条误识别，原因为查询文本同时包含"IRR"和"总投资"两个指标词，模板匹配选择优先级冲突。}
	\end{center}
\end{table}

\subsection{\heiti 案例分析：IRR路径追溯}

以IRR（内部收益率）指标为例，演示完整的计算路径追溯。IRR位于表5-利润表-资本金的D105单元格，其正向追溯路径覆盖了从收入、成本到利润的完整计算链：

\begin{enumerate}
	\item IRR(D105) $\leftarrow$ 各期净现金流(D52-D104)
	\item 各期净现金流 $\leftarrow$ 净利润 + 折旧摊销 $-$ 资本支出
	\item 净利润 $\leftarrow$ 营业收入 $-$ 总成本 $-$ 税金
	\item 营业收入 $\leftarrow$ 电价 $\times$ 售电量
	\item 总成本 $\leftarrow$ 运维成本 + 财务费用 + 折旧摊销
\end{enumerate}

通过知识图谱的 \textit{depends\_on} 关系多跳查询，IRR完整计算链可在<0.5秒内返回。相比之下，人工追溯需依次打开14个工作表、定位相关公式、逐个确认依赖关系，耗时约45-60分钟。效率提升超5,000倍。

反向影响分析同样高效：修改参数输入表中的某利率参数，系统沿 \textit{depends\_on} 关系正向追溯，1.2秒内返回该参数影响的1,021个下游单元格，覆盖资金筹措表、折旧摊销表、成本费用表等5个工作表。


\section{\heiti 结论}

本文选择上下文规则优先的混合分类策略，其根本依据在于：能源财务模型的公式业务含义由行列标题决定而非由公式语法决定——消融实验中同一SUM函数在不同行代表完全不同的业务类型，这一结构性特征使得纯语法方法从根本上不适用于本场景，而非仅仅在准确率上略有不足。

消融实验揭示的核心规律是：跨业务Sheet现象（单一Sheet内嵌多个业务模块）构成了规则方法的天然边界。本模型时间序列Sheet内嵌折旧时间轴逻辑，导致28/35条错误集中于此。这一发现说明，知识图谱构建中的语义歧义问题往往来自模型设计层面的约定，而非方法本身的局限——这对财务模型的规范化设计具有直接的实践指导意义。

知识图谱的工程价值不仅体现在查询响应时间（IRR路径追溯<0.5秒）上，更体现在使财务模型从单点专家知识转变为可查询、可审计的组织知识资产——这是传统Excel文件无法实现的根本转变。当前模型包含289个含错误公式，这些错误在人工审计中极难全面检出，而知识图谱的依赖图可以直接定位错误节点对下游指标的影响范围，这一能力在本文实验中尚未充分展开，是后续研究的重点方向之一。

本研究的局限在于两点：其一，研究对象为单一抽水蓄能项目模型，风电、光伏、核电等类型的财务模型是否存在相同的跨业务Sheet现象尚待验证；其二，LLM增强分类模块当前仅有框架设计，21.7\%的高歧义公式的实际分类效果有待实验验证。后续工作将重点收集多类型能源财务模型，构建多案例对比基准，同时探索基于LLM的跨业务公式语义消歧方法。

\vspace {3mm}
\zihao{5}{
	\noindent \textsf{致\quad 谢}\quad \textit{感谢提供真实抽水蓄能电站财务评价模型数据的研究合作方，以及在数据标注和实验验证过程中给予支持的团队成员。}}

\vspace {5mm}
\centerline
{\zihao{5}\textsf{参~考~文~献}}
\zihao{5-} \addtolength{\itemsep}{-1em}
\vspace {1.5mm}
\bibliographystyle{gbt7714-numerical}
\bibliography{ref.bib}


\noindent {\zihao{5}\bf{附录 A}.}

{\zihao{5-}\setlength\parindent{2em}
	\textbf{消融实验完整结果（n=299，按Sheet分层）}

	表~\ref{tab:appendix} 展示了消融实验在各Sheet上的分层结果。时间序列Sheet是错误集中区域，28/35条方法B/C错误均源于"跨业务Sheet"现象——该Sheet内嵌折旧摊销时间轴计算逻辑，导致Sheet名称与行标题语义冲突。

	\begin{table}[H]
		\centering
		\caption{消融实验分层结果}
		\label{tab:appendix}
		\vspace {-2.5mm}
		\begin{center}
			\begin{tabular}{lccc}
				\toprule
				Sheet名称 & 方法A/\% & 方法B/C/\% & 错误数 \\
				\hline
				参数输入表 & 25.0 & 91.7 & 1 \\
				时间序列 & 12.3 & 72.9 & 28 \\
				投产\&达产比例 & 18.5 & 90.2 & 3 \\
				表1-资金筹措及还本付息表 & 22.1 & 91.0 & 2 \\
				表2-折旧摊销表 & 19.8 & 89.5 & 4 \\
				表3-成本费用表 & 20.5 & 90.1 & 3 \\
				表4-收入税金表 & 21.0 & 91.3 & 1 \\
				表5-利润表-资本金 & 24.3 & 88.9 & 2 \\
				表6-现金流量表-资本金 & 23.7 & 87.5 & 2 \\
				\bottomrule
			\end{tabular}
		\end{center}
	\end{table}
}

%\begin{biography}[yourphotofilename.jpg]
%	\noindent
%	\textbf{Author 1}\ \ (Anonymized for review.)
%\end{biography}
%
%\begin{biography}[yourphotofilename.jpg]
%	\noindent
%	\textbf{Author 2}\ \ (Anonymized for review.)
%\end{biography}
%
%\begin{biography}[yourphotofilename.jpg]
%	\noindent
%	\textbf{Author 3}\ \ (Anonymized for review.)
%\end{biography}

\zihao{5}
\noindent \textbf{Background}

\zihao{5-}{
	\setlength\parindent{2em}
	This paper addresses the problem of knowledge extraction and semantic tracing for Excel-based financial models in the energy sector. Excel spreadsheets have long been the de facto standard for financial evaluation in energy projects, but their potential as structured knowledge sources has not been fully exploited.

	With China's pumped-storage hydropower installed capacity exceeding 50 million kilowatts, the scale and complexity of financial evaluation models have grown significantly. The model studied in this paper comprises 42,109 formulas and 9,199 cross-sheet references, requiring 2--3 person-days for manual audit. Existing spreadsheet analysis methods focus on formula correctness and lack semantic understanding capabilities, resulting in a severe knowledge extraction gap.

	This paper proposes a knowledge graph-driven approach that transforms Excel financial models into searchable, auditable, and traceable knowledge assets. The method integrates DAG-based dependency parsing, context-rule-priority semantic classification, and knowledge graph construction. Experimental results demonstrate that the proposed method achieves 88.3\% semantic classification accuracy (vs. 20.7\% for pure formula-structure methods), constructs a knowledge graph of 42,138 entities and 65,333 relations, and enables sub-second IRR computation path tracing (over 5,000$\times$ faster than manual tracing), providing a systematic solution for building an intelligent financial model audit loop.}

\end{document}
